\documentclass[11pt]{article}

% basic packages
\usepackage[margin=1in]{geometry}
\usepackage[pdftex]{graphicx}
\usepackage{amsmath,amssymb,amsthm}
\usepackage{booktabs,longtable}
\usepackage{custom}

% page formatting
\usepackage{fancyhdr}
\pagestyle{fancy}

\renewcommand{\sectionmark}[1]{\markright{\textsf{\thesection. #1}}}
\renewcommand{\subsectionmark}[1]{}
\lhead{\textbf{\thepage} \ \ \nouppercase{\rightmark}}
\chead{}
\rhead{}
\lfoot{}
\cfoot{}
\rfoot{}
\setlength{\headheight}{14pt}

\linespread{1.00}
\setlength{\parindent}{0.2in} % reduce paragraph indent a bit
\setlength{\LTpre}{0.4em}
\setlength{\LTpost}{0.4em}
\setcounter{secnumdepth}{2} % no numbered subsubsections
\setcounter{tocdepth}{2} % no subsubsections in ToC

\begin{document}

% make title page
\thispagestyle{empty}
\bigskip
\vspace{0.1cm}

{\fontsize{18}{18} \selectfont \bf \sffamily Model Equations}

% main content
\section*{Common conventions and units \cite{BishopVanHorn2023}}

All three multicomponent models are implemented in dimensionless form.  Throughout
this document,
\begin{equation}
    D_i \equiv \partial_i-iA_i,\qquad
    D_t \equiv \partial_t+i\mu,\qquad i\in\{x,y\},
    \label{eq:covariant-derivatives}
\end{equation}
where $\mu$ is the scalar electric potential (the same quantity is sometimes
denoted $\phi$).  Complex conjugation is denoted by an asterisk.  The gauge
transformation
\begin{equation}
    \psi_n\rightarrow e^{i\chi}\psi_n,\qquad
    \mathbf A\rightarrow\mathbf A+\boldsymbol\nabla\chi,\qquad
    \mu\rightarrow\mu-\partial_t\chi
\end{equation}
leaves all of the equations below invariant.

The pyTDGL dimensional scales, in SI, are~\cite{BishopVanHorn2023}
\begin{align}
    \tau_0 &= \mu_0\sigma\lambda^2,
    &B_0=B_{c2}&=\frac{\Phi_0}{2\pi\xi^2},
    &A_0&=\xi B_{c2}=\frac{\Phi_0}{2\pi\xi},\\
    J_0&=\frac{4\xi B_{c2}}{\mu_0\lambda^2},
    &K_0&=J_0d=\frac{4\xi B_{c2}}{\mu_0\Lambda},
    &V_0&=\frac{\xi J_0}{\sigma}
        =\frac{4\xi^2B_{c2}}{\mu_0\sigma\lambda^2},
    \label{eq:physical-scales}
\end{align}
where $\xi$ is the coherence length, $\lambda$ the London penetration depth,
$d$ the film thickness, $\Lambda=\lambda^2/d$, $\sigma$ the normal-state
conductivity, and $\Phi_0=h/(2e)$ the flux quantum.  Solver variables are
\begin{equation}
    \mathbf r=\frac{\mathbf r_{\rm phys}}{\xi},\quad
    t=\frac{t_{\rm phys}}{\tau_0},\quad
    \mathbf A=\frac{\mathbf A_{\rm phys}}{A_0},\quad
    \mathbf B=\frac{\mathbf B_{\rm phys}}{B_0},\quad
    \mu=\frac{V_{\rm phys}}{V_0},\quad
    \mathbf j=\frac{\mathbf K_{\rm phys}}{K_0}.
    \label{eq:dimensionless-variables}
\end{equation}
Order parameters and code coefficients are dimensionless.  The source papers
use Gaussian cgs; all dimensional formulas below are converted to SI, while
dimensionless ratios are unit-system independent.

For each model, current conservation is imposed through
\begin{align}
    \mathbf j_n&=-\boldsymbol\nabla\mu-\partial_t\mathbf A,\\
    \boldsymbol\nabla^2\mu
        &=\boldsymbol\nabla\!\cdot\mathbf j_s
          -\boldsymbol\nabla\!\cdot\partial_t\mathbf A,
    \label{eq:poisson-common}
\end{align}
where $\mathbf j_s$ is the transport-normalized supercurrent specified for
each model below.  The multicomponent solvers presently use a prescribed
vector potential; self-consistent magnetic screening is not implemented.

\section{s+d \cite{Goncalves2014,Zhu1999}}

\subsection{Order-parameter equations}

Let $\psi_d$ and $\psi_s$ denote the dimensionless $d$- and $s$-wave order
parameters.  The implemented mixed-symmetry TDGL equations follow the
dimensionless $d+s$ framework of Refs.~\cite{Goncalves2014,Zhu1999}:
\begin{align}
    D_t\psi_d={}&(D_x^2+D_y^2)\psi_d
      +\varepsilon\psi_d-|\psi_d|^2\psi_d
      -\frac{\tau_3}{2}|\psi_s|^2\psi_d
      -\tau_4\psi_s^2\psi_d^*
      +\eta_v(D_x^2-D_y^2)\psi_s,
    \label{eq:sd-d}\\
    \eta_sD_t\psi_s={}&\eta_s(D_x^2+D_y^2)\psi_s
      +\nu\psi_s-\tau_1|\psi_s|^2\psi_s
      -\frac{\tau_3}{2}|\psi_d|^2\psi_s
      -\tau_4\psi_d^2\psi_s^*
      +\eta_v(D_x^2-D_y^2)\psi_d.
    \label{eq:sd-s}
\end{align}
Here $\varepsilon=\varepsilon(\mathbf r,t)$ is the solver input
\texttt{disorder\_epsilon}; it changes only the quadratic drive of the
$d$ sector.  The standard homogeneous value is $\varepsilon=1$.

Define the mixed directional currents
\begin{equation}
    Q_i=\operatorname{Im}\!\left(
        \psi_d^*D_i\psi_s+\psi_s^*D_i\psi_d
    \right).
\end{equation}
The condensate current used by the implementation before electromagnetic
normalization is
\begin{align}
    J_{{\rm cond},x}&=
      \operatorname{Im}(\psi_d^*D_x\psi_d)
      +\eta_s\operatorname{Im}(\psi_s^*D_x\psi_s)-\eta_vQ_x,\\
    J_{{\rm cond},y}&=
      \operatorname{Im}(\psi_d^*D_y\psi_d)
      +\eta_s\operatorname{Im}(\psi_s^*D_y\psi_s)+\eta_vQ_y.
\end{align}
The transport-normalized current stored by the solver and used in
Eq.~\eqref{eq:poisson-common} is
\begin{equation}
    \mathbf j_s=\frac{\mathbf J_{\rm cond}}{\beta_{\rm em}}.
    \label{eq:sd-current}
\end{equation}
The opposite signs of the $x$ and $y$ mixed-current terms encode the
$d_{x^2-y^2}$ symmetry.

\subsection{Dimensional origin and SI units}

From effective masses $m_d^*,m_s^*,m_v^*$, quadratic coefficients
$\alpha_d,\alpha_s$, quartic coefficients $\beta_1,\ldots,\beta_4$, diffusion
coefficient $\mathcal D$, and conductivity $\sigma$, Gon\c{c}alves
\textit{et al.} obtain~\cite{Goncalves2014}
\begin{align}
    \eta_s&=\frac{m_d^*}{m_s^*}, &
    \eta_v&=\frac{m_d^*}{m_v^*}, &
    \nu&=\frac{\alpha_s}{\alpha_d},\\
    \tau_1&=\frac{\beta_1}{\beta_2}, &
    \tau_3&=\frac{\beta_3}{\beta_2}, &
    \tau_4&=\frac{\beta_4}{\beta_2},
    \label{eq:sd-parameter-map}\\
    \beta
      &=\mu_0\sigma\kappa_d^2\mathcal D
       =\frac{\mu_0\sigma\lambda_d^2}
              {\xi_d^2/\mathcal D}.
    \label{eq:sd-beta-map}
\end{align}
$\beta$ is the ratio of the electromagnetic relaxation time
$\mu_0\sigma\lambda_d^2$ to the order-parameter diffusion time
$\xi_d^2/\mathcal D$; Eq.~\eqref{eq:sd-beta-map} is the SI conversion of the
paper's expression.

The remaining source scales are
\begin{align}
    |\psi_{d0}|^2&=\frac{\alpha_d}{2\beta_2}, &
    \xi_d^2&=\frac{\hbar^2}{2m_d^*\alpha_d}, &
    \lambda_d^2&=\frac{2m_d^*\beta_2}
        {\mu_0(e^*)^2\alpha_d},\\
    \kappa_d&=\frac{\lambda_d}{\xi_d}, &
    B_{c2}&=\frac{\Phi_0}{2\pi\xi_d^2}, &
    \Phi_0&=\frac{h}{2e}.
    \label{eq:sd-source-scales}
\end{align}
SI units are: $m_i^*$ [kg], $\mathcal D$ [$\mathrm{m^2\,s^{-1}}$], $\sigma$
[$\mathrm{S\,m^{-1}}$], $\xi_d,\lambda_d$ [m], $B$ [T], $\mathbf A$
[$\mathrm{T\,m}$], and $\Phi_0$ [Wb].  Each set
$(\alpha_s,\alpha_d)$ and $(\beta_1,\ldots,\beta_4)$ must use common compatible
units; their ratios are dimensionless for any consistent order-parameter
normalization.

\subsection{Parameters}

N.B. Theoretical values for the parameters have been derived microscopically, corresponding to those in section VIII of Goncalves, but we allow them to be adjusted just in case.

{\small
\begin{longtable}{@{}p{0.22\textwidth}p{0.59\textwidth}p{0.11\textwidth}@{}}
\toprule
Code name & Meaning & Default\\
\midrule
\endhead
$\eta_s$ (\texttt{eta\_s}) & Effective-mass ratio $m_d^*/m_s^*$; sets relative $s$-sector stiffness and the coefficient of $D_t\psi_s$. & $1$\\
$\eta_v$ (\texttt{eta\_v}) & Mixed-gradient effective-mass ratio $m_d^*/m_v^*$, responsible for the anisotropic $s$--$d$ coupling. & $0$\\
$\nu$ (\texttt{nu}) & Quadratic ratio $\alpha_s/\alpha_d$; gives the free-energy term $-\nu|\psi_s|^2$. & $-1$\\
$\tau_1$ (\texttt{tau1}) & Quartic ratio $\beta_1/\beta_2$, comparing the $s$ self-interaction with the $d$ self-interaction. & $1$\\
$\tau_3$ (\texttt{tau3}) & Density-coupling ratio $\beta_3/\beta_2$; TDGL cross-density coefficient $\tau_3/2$. & $0$\\
$\tau_4$ (\texttt{tau4}) & Phase-sensitive quartic ratio $\beta_4/\beta_2$, coupling $\psi_s^2\psi_d^*$ and $\psi_d^2\psi_s^*$. & $0$\\
$\beta_{\rm em}$ (\texttt{beta\_em}) & SI electromagnetic factor $\beta=\mu_0\sigma\kappa_d^2\mathcal D$.  Since the solver prescribes $\mathbf A$ rather than evolving the paper's $\beta\partial_t\mathbf A$ equation, this field acts only as the current divisor in Eq.~\eqref{eq:sd-current}. & $1$\\
$\varepsilon$ (\texttt{disorder\_\allowbreak epsilon}) & Local $d$-sector quadratic coefficient, interpretable as $\alpha_d(\mathbf r,t)/\alpha_{d,\rm ref}$; the homogeneous Gon\c{c}alves normalization gives $\varepsilon=1$. & $1$\\
\bottomrule
\end{longtable}
}

The validation conditions are
\begin{equation}
    \eta_s>0,\qquad \eta_s>\eta_v^2,\qquad
    \tau_1>0,\qquad \beta_{\rm em}>0,
\end{equation}
with $\varepsilon\leq1$ enforced by the solver.  The condition
$\eta_s>\eta_v^2$ makes the mixed-gradient sector positive definite.

\subsection{Limitations and approximations}

The Gon\c{c}alves formulation treats the diffusion coefficient and conductivity
phenomenologically and does not impose a strict clean- or dirty-limit condition.
The microscopic theory of Zhu \textit{et al.} includes magnetic and nonmagnetic
impurity scattering rather than assuming only the asymptotic dirty limit.
Nevertheless, the implementation's local Ohmic normal current and diffusive
relaxation are dirty-limit-like; quantitative dynamics require material-specific
scattering and relaxation parameters.

The field $\varepsilon(\mathbf r,t)$ changes only the $d$-sector quadratic
coefficient and is not a complete microscopic disorder model.  The factor
$\beta_{\rm em}$ is only the current normalization in
Eq.~\eqref{eq:sd-current}.  A scalar terminal value fixes $\psi_d$ and sets
$\psi_s=0$ at terminal sites, which is a numerical contact model rather than a
general $s+d$ interface condition.

Note that $H_{c2}$ here is for a pure d-wave and therefore the ``actual" $H_{c2}$ could be different (a possible simulation result).

\section{d+d' \cite{Lei2000, Wang1999Orbital}}

\subsection{Free energy and equilibrium-relaxation equations}

Let $d$ denote the dominant $d_{x^2-y^2}$ component and $d'$ the
subdominant $d_{xy}$ component.  In the code they are stored as
$\psi_1=d$ and $\psi_2=d'$, respectively.  The dimensionless free-energy
density of Lei, Aruna, and Wang~\cite{Lei2000} is augmented by the optional
orbital-Zeeman coupling derived microscopically by Wang and
Wang~\cite{Wang1999Orbital}:
\begin{align}
    f_{d+d'}={}&-|d|^2-\alpha|d'|^2
      +\frac{1}{2}\left(|d|^4+|d'|^4\right)
      +\frac{1}{3}|d|^2|d'|^2
      +\frac{1}{6}\left(d^*d'+d d'^*\right)^2 \notag\\
      &+|\mathbf Dd|^2+|\mathbf Dd'|^2+\kappa^2b^2
      -i g_Zb\left(d^*d'-d d'^*\right),
    \label{eq:ddp-free-energy}
\end{align}
where
\begin{equation}
    b=(\boldsymbol\nabla\times\mathbf A)_z
      =\frac{B_z}{B_{c2}},
    \qquad B_{c2}=\frac{\Phi_0}{2\pi\xi^2}.
    \label{eq:ddp-field}
\end{equation}
The original field-driven model in Ref.~\cite{Lei2000} deliberately omits the
last term; setting $g_Z=0$ reproduces that functional.  The signed coefficient
$g_Z$ incorporates the microscopic particle--hole asymmetry and the
dimensionless normalization of the orbital-Zeeman interaction.  The code
computes $b$ at mesh sites by a discrete Stokes curl of the edge-centered
applied vector potential.

The source papers specify a static GL functional rather than physical TDGL
kinetics.  The implemented equations are therefore phenomenological negative
gradient flow for finding stationary states:
\begin{align}
    r_dD_td={}&(D_x^2+D_y^2)d+d-|d|^2d
      -\frac{2}{3}|d'|^2d-\frac{1}{3}d'^2d^*
      +i g_Zb d',
    \label{eq:ddp-d}\\
    r_{d'}D_td'={}&(D_x^2+D_y^2)d'+\alpha d'-|d'|^2d'
      -\frac{2}{3}|d|^2d'-\frac{1}{3}d^2d'^*
      -i g_Zb d.
    \label{eq:ddp-dprime}
\end{align}
The factors $2/3$ and $1/3$ follow by differentiating both quartic cross
terms in Eq.~\eqref{eq:ddp-free-energy}; they are not independent fitting
parameters.

At $b=0$ and $g_Z=0$, a uniform pure-$d$ stationary state has
$|d|=1$ and $d'=0$.  For $1/3<\alpha<1$, the uniform mixed minimum instead
has
\begin{equation}
    |d|^2=\frac{3(3-\alpha)}{8},\qquad
    |d'|^2=\frac{3(3\alpha-1)}{8},\qquad
    \arg(d')-\arg(d)=\pm\frac{\pi}{2}.
    \label{eq:ddp-uniform-minimum}
\end{equation}
For $g_Zb>0$, the Zeeman term lifts this chiral degeneracy and favors
$d'/d=-i|d'/d|$; reversing the field reverses the favored chirality.

\subsection{Current and orbital magnetization}

Because Eq.~\eqref{eq:ddp-free-energy} has equal diagonal stiffnesses and no
mixed-gradient term, the implemented condensate and transport currents are
\begin{align}
    \mathbf J_{\rm cond}
      &=\operatorname{Im}(d^*\mathbf Dd)
        +\operatorname{Im}(d'^*\mathbf Dd'),\\
    \mathbf j_s&=\frac{\mathbf J_{\rm cond}}{c_{\rm em}},
    \qquad c_{\rm em}\equiv\texttt{em\_coupling}.
    \label{eq:ddp-current}
\end{align}
The local dimensionless orbital magnetization associated with the Zeeman term
is
\begin{equation}
    m_z=-\frac{\partial f_Z}{\partial b}\bigg|_{d,d'}
       =i g_Z\left(d^*d'-d d'^*\right),
    \qquad f_Z=-i g_Zb\left(d^*d'-d d'^*\right),
    \label{eq:ddp-magnetization}
\end{equation}
This is the order-parameter orbital contribution and excludes the derivative
of the separately prescribed $\kappa^2b^2$ field energy.
Equation~\eqref{eq:ddp-magnetization} is exposed as the diagnostic
\texttt{Solution.\allowbreak orbital\_\allowbreak magnetization}.

The discrete condensate free-energy diagnostic includes the local potential,
gradient energy, and orbital-Zeeman energy in
Eq.~\eqref{eq:ddp-free-energy}.  It omits $\kappa^2b^2$, which is constant with
respect to $d,d'$ when $\mathbf A$ is prescribed and therefore does not affect
order-parameter relaxation.

\subsection{Parameters}

{\small
\begin{longtable}{@{}p{0.27\textwidth}p{0.55\textwidth}p{0.10\textwidth}@{}}
\toprule
Code name & Meaning & Default\\
\midrule
\endhead
$\alpha$ (\texttt{alpha}) & Ratio of the subdominant and dominant quadratic GL coefficients in the normalization of Ref.~\cite{Lei2000}; controls the zero-field $d_{xy}$ instability. & $0.5$\\
$r_d$ (\texttt{relaxation\_d}) & Positive phenomenological relaxation coefficient for the dominant $d_{x^2-y^2}$ component. & $1$\\
$r_{d'}$ (\texttt{relaxation\_d\_prime}) & Positive phenomenological relaxation coefficient for the subdominant $d_{xy}$ component. & $1$\\
$c_{\rm em}$ (\texttt{em\_coupling}) & Implementation-specific conversion from condensate-current normalization to the transport normalization used by the Poisson solver. & $1$\\
$g_Z$ (\texttt{zeeman\_coupling}) & Signed dimensionless orbital-Zeeman coefficient; zero selects the unmodified Lei--Aruna--Wang model. & $0$\\
\bottomrule
\end{longtable}
}

All parameters must be finite, $r_d,r_{d'},c_{\rm em}>0$, and the current
implementation requires $|g_Z|<1$.  Negative $\alpha$ is permitted.  Values
$\alpha>1$ are numerically allowed with a warning because they leave the
assumed subdominant-$d_{xy}$ regime of the source model.

\subsection{Limitations and approximations}

Our equations are weak-coupling microscopic results based on BCS, Gor'kov, and path-integral
methods without impurity scattering; they do not assume the dirty limit.  The
coefficients $r_d,r_{d'}$ are phenomenological additions, so the simulated
trajectory is equilibrium relaxation rather than a microscopic clean-limit
TDGL prediction.

The Zeeman term is optional because it changes the phase diagram qualitatively:
when $g_Zb\ne0$, it acts as a linear source for the opposite component and
removes the exact degeneracy of the two chiralities.  Comparisons with the
phase boundary of Ref.~\cite{Lei2000} must use $g_Z=0$; nonzero values instead
implement the microscopic extension of Ref.~\cite{Wang1999Orbital}.

The model has equal component stiffnesses, no mixed-gradient term, and fixed
quadratic coefficients, so nonuniform \texttt{disorder\_epsilon} is unsupported.
The reported orbital magnetization omits its bound current from the transport
Poisson solve.  Transport terminals require \texttt{terminal\_psi=None} and
natural order-parameter boundary conditions.

\section{s+s \cite{Zhitomirsky2004,Kogan2011}}

\subsection{Free energy and order-parameter equations}

Let $\psi_1$ and $\psi_2$ denote the two dimensionless isotropic $s$-wave
band order parameters.  With $k_1=1$ and
$k_2=k_{2/1}\equiv\texttt{k2\_over\_k1}$, the dimensionless condensate
free-energy density follows Refs.~\cite{Zhitomirsky2004,Kogan2011}:
\begin{equation}
    f=\sum_{n=1}^{2}\left[
        a_n|\psi_n|^2+\frac{b_n}{2}|\psi_n|^4
        +k_n|\mathbf D\psi_n|^2
    \right]
    -\gamma_{12}\left(\psi_1\psi_2^*+\psi_2\psi_1^*\right).
    \label{eq:ss-free-energy}
\end{equation}
The implemented dissipative extension is the negative gradient flow of
Eq.~\eqref{eq:ss-free-energy}:
\begin{align}
    r_1D_t\psi_1={}&(D_x^2+D_y^2)\psi_1
      -a_1\psi_1-b_1|\psi_1|^2\psi_1+\gamma_{12}\psi_2,
    \label{eq:ss-1}\\
    r_2D_t\psi_2={}&k_{2/1}(D_x^2+D_y^2)\psi_2
      -a_2\psi_2-b_2|\psi_2|^2\psi_2+\gamma_{12}\psi_1.
    \label{eq:ss-2}
\end{align}
The unscaled condensate current and the transport-normalized current are
\begin{align}
    \mathbf J_{\rm cond}
      &=\operatorname{Im}(\psi_1^*\mathbf D\psi_1)
        +k_{2/1}\operatorname{Im}(\psi_2^*\mathbf D\psi_2),\\
    \mathbf j_s&=\frac{\mathbf J_{\rm cond}}{c_{\rm em}},
    \qquad c_{\rm em}\equiv\texttt{em\_coupling}.
    \label{eq:ss-current}
\end{align}
A positive $\gamma_{12}$ minimizes the energy for equal component phases; a
negative value favors a relative phase of $\pi$.

\subsection{Dimensional origin and SI units}

The Kogan--Schmalian two-band functional, converted from Gaussian cgs to SI,
is
\begin{equation}
    F_{\rm SI}=\int dV\left\{
      \sum_{\nu=1}^{2}\left[
        a_\nu^{(G)}|\Delta_\nu|^2
        +\frac{b_\nu^{(G)}}{2}|\Delta_\nu|^4
        +K_\nu^{(G)}|\boldsymbol\Pi\Delta_\nu|^2
      \right]
      -\gamma^{(G)}(\Delta_1\Delta_2^*+\Delta_2\Delta_1^*)
      +\frac{B^2}{2\mu_0}
    \right\},
    \label{eq:ss-dimensional-si}
\end{equation}
where $\boldsymbol\Pi=\boldsymbol\nabla+2\pi i\mathbf A/\Phi_0$,
$\Phi_0=h/(2e)$, and $\Delta_\nu$ is the dimensional band gap.  Superscript
$(G)$ labels dimensional GL coefficients, not powers.

For the clean, isotropic weak-coupling model, the paper gives
~\cite{Kogan2011}
\begin{align}
    a_\nu^{(G)}
      &=\frac{N(0)}{\eta_\lambda}
        \left(S_\nu-n_\nu\eta_\lambda\tau\right),&
    b_\nu^{(G)}&=\frac{N(0)}{W^2}n_\nu,\\
    \gamma^{(G)}&=\frac{N(0)}{\eta_\lambda}\lambda_{12},&
    K_\nu^{(G)}&=\frac{N(0)\hbar^2v_\nu^2}{6W^2}n_\nu,
    \label{eq:ss-microscopic-coefficients}\\
    \eta_\lambda&=\lambda_{11}\lambda_{22}-\lambda_{12}^2,&
    S_1&=\lambda_{22}-n_1\eta_\lambda S,\qquad
    S_2=\lambda_{11}-n_2\eta_\lambda S,\\
    W^2&=\frac{8\pi^2(k_BT_c)^2}{7\zeta(3)},&
    S&=\ln\!\left(\frac{2e^{\gamma_E}\hbar\omega_D}
                         {\pi k_BT_c}\right),\qquad
    \tau=1-\frac{T}{T_c}.
\end{align}
Here $N(0)$ is the total single-spin density of states at the Fermi level;
$n_\nu=N_\nu/N(0)$ the band fraction; $\lambda_{\nu\mu}$ the dimensionless
pairing matrix; $v_\nu$ the Fermi velocity; $\omega_D$ the cutoff frequency;
and $\gamma_E$ Euler's constant.  Equations~\eqref{eq:ss-microscopic-coefficients}
are one weak-coupling realization; Eq.~\eqref{eq:ss-dimensional-si} also
applies phenomenologically.

SI units are: $\Delta_\nu,W,k_BT_c$ [J], $N(0)$
[$\mathrm{J^{-1}m^{-3}}$], $a_\nu^{(G)},\gamma^{(G)}$
[$\mathrm{J^{-1}m^{-3}}$], $b_\nu^{(G)}$
[$\mathrm{J^{-3}m^{-3}}$], $K_\nu^{(G)}$
[$\mathrm{J^{-1}m^{-1}}$], $v_\nu$ [$\mathrm{m\,s^{-1}}$], $\mathbf A$
[$\mathrm{T\,m}$], $B$ [T], and $F_{\rm SI}$ [J].
$n_\nu,\lambda_{\nu\mu},\eta_\lambda,S,\tau$ are dimensionless.

With gap scale $\Delta_0$, length $\xi$, and energy-density scale
$f_0=K_1^{(G)}\Delta_0^2/\xi^2$, the code coefficients are
\begin{align}
    a_\nu&=\frac{a_\nu^{(G)}\xi^2}{K_1^{(G)}},&
    b_\nu&=\frac{b_\nu^{(G)}\Delta_0^2\xi^2}{K_1^{(G)}},\\
    k_{2/1}&=\frac{K_2^{(G)}}{K_1^{(G)}},&
    \gamma_{12}&=\frac{\gamma^{(G)}\xi^2}{K_1^{(G)}}.
    \label{eq:ss-code-map}
\end{align}
Other scales are equivalent if the whole functional is rescaled consistently.

\subsection{Parameters}

{\small
\begin{longtable}{@{}p{0.27\textwidth}p{0.55\textwidth}p{0.10\textwidth}@{}}
\toprule
Code name & Meaning & Default\\
\midrule
\endhead
$a_1$ (\texttt{a1}), $a_2$ (\texttt{a2}) & Dimensionless $a_\nu^{(G)}$; contain a constant interband contribution and a term proportional to $1-T/T_c$. & $-1,-1$\\
$b_1$ (\texttt{b1}), $b_2$ (\texttt{b2}) & Dimensionless quartic self-interactions corresponding to $b_\nu^{(G)}=N(0)n_\nu/W^2$. & $1,1$\\
$k_{2/1}$ (\texttt{k2\_over\_k1}) & Stiffness ratio $K_2^{(G)}/K_1^{(G)}$; depends on band fractions and squared Fermi velocities. & $1$\\
$\gamma_{12}$ (\texttt{josephson\_gamma}) & Dimensionless interband Josephson coupling corresponding to $\gamma^{(G)}=N(0)\lambda_{12}/\eta_\lambda$. & $0$\\
$r_1$ (\texttt{relaxation1}), $r_2$ (\texttt{relaxation2}) & Phenomenological TDGL kinetic coefficients; absent from the static Kogan--Schmalian functional. & $1,1$\\
$c_{\rm em}$ (\texttt{em\_coupling}) & Implementation-specific conversion from condensate-current normalization to the transport normalization used by the Poisson solver; it is not a coefficient in the static source functional. & $1$\\
\bottomrule
\end{longtable}
}

\enlargethispage{3\baselineskip}
The validation conditions are
\begin{equation}
    b_1>0,\qquad b_2>0,\qquad k_{2/1}>0,\qquad
    r_1>0,\qquad r_2>0,\qquad c_{\rm em}>0.
\end{equation}
$a_1,a_2$ may have either sign; an uncoupled uniform band with $a_n<0$ has
$|\psi_n|=\sqrt{-a_n/b_n}$.

\subsection{Limitations and approximations}

The microscopic coefficient map in Ref.~\cite{Kogan2011} is derived for a
clean, isotropic, weak-coupling material, not the dirty limit.  The GL result
that the two gaps are proportional near the common $T_c$ is more general and
also has dirty-limit derivations, but those dirty microscopic coefficients are
not encoded here.  The relaxation coefficients $r_1,r_2$ remain
phenomenological.

The fixed-temperature coefficients $a_1,a_2$ require
\texttt{disorder\_epsilon}$=1$, and the model omits interband gradient drag and
impurity-induced interband scattering.  Near $T_c$, the two stored fields do
not imply two independent GL coherence lengths.  The factor $c_{\rm em}$ is a
current normalization, and scalar terminal order-parameter values are
unsupported because their two-band mapping is unspecified.

\begin{thebibliography}{99}

\bibitem{BishopVanHorn2023}
L.~Bishop-Van Horn,
``pyTDGL: Time-dependent Ginzburg--Landau in Python,''
\textit{Computer Physics Communications} \textbf{291}, 108799 (2023).
\href{https://doi.org/10.1016/j.cpc.2023.108799}{doi:10.1016/j.cpc.2023.108799}.

\bibitem{Goncalves2014}
W.~C. Gon\c{c}alves, E.~Sardella, V.~F. Becerra, M.~V. Milosevic,
and F.~M. Peeters,
``Numerical solution of the time dependent Ginzburg--Landau equations for
mixed ($d+s$)-wave superconductors,''
\textit{Journal of Mathematical Physics} \textbf{55}, 041501 (2014).
\href{https://doi.org/10.1063/1.4870874}{doi:10.1063/1.4870874}.

\bibitem{Zhu1999}
J.-X. Zhu, W.~Kim, C.~S. Ting, and C.-R. Hu,
``Time-dependent Ginzburg--Landau equations for mixed $d$- and $s$-wave
superconductors,''
\textit{Physical Review B} \textbf{59}, 11527--11534 (1999).
\href{https://doi.org/10.1103/PhysRevB.59.11527}{doi:10.1103/PhysRevB.59.11527}.

\bibitem{Lei2000}
B.~Lei, S.~A. Aruna, and Q.-H. Wang,
``Field driven pairing state phase transition in
$d_{x^2-y^2}+id_{xy}$-wave superconductors,''
\textit{Physical Review B} \textbf{62}, 8687 (2000).
\href{https://doi.org/10.1103/PhysRevB.62.8687}{doi:10.1103/PhysRevB.62.8687}.

\bibitem{Wang1999Orbital}
Q.-H. Wang and Z.~D. Wang,
``Zeeman coupling and abnormal thermal conductivities in BSCCO
superconductors,''
arXiv:cond-mat/9909399 (1999).
\href{https://doi.org/10.48550/arXiv.cond-mat/9909399}{doi:10.48550/arXiv.cond-mat/9909399}.

\bibitem{Zhitomirsky2004}
M.~E. Zhitomirsky and V.-H. Dao,
``Ginzburg--Landau theory of vortices in a multigap superconductor,''
\textit{Physical Review B} \textbf{69}, 054508 (2004).
\href{https://doi.org/10.1103/PhysRevB.69.054508}{doi:10.1103/PhysRevB.69.054508}.

\bibitem{Kogan2011}
V.~G. Kogan and J.~Schmalian,
``Ginzburg--Landau theory of two-band superconductors,''
\textit{Physical Review B} \textbf{83}, 054515 (2011).
\href{https://doi.org/10.1103/PhysRevB.83.054515}{doi:10.1103/PhysRevB.83.054515}.

\end{thebibliography}

\end{document}
